\chapter{Bevezetés}\label{ch:introduction}

A szakdolgozatom célja egy könnyen kezelhető receptkezelő webes alkalmazás létrehozása. Az alkalmazás lehetővé teszi a felhasználók számára, hogy saját receptjeiket feltöltsék, megosszák másokkal, valamint értékeljék és kommentálják a meglévő recepteket. Emellett az alkalmazás támogatja a keresést és szűrést, különösen nagy hangsúlyt fektetve az allergén információk megjelenítésére, hogy az ételallergiával élő felhasználók is biztonságosan használhassák a platformot.

\section{Motiváció}

Családunkban több generáció házi, sokszor kézzel írt receptjeit használja. Ezek a receptek gyakran papíron vagy különböző digitális formátumokban vannak tárolva, melyek követése és megosztása nehézkes lehet. Ezek utólagos módosítása szintén nehézkesen, vagy nehezen követhető módon történik.
A jelenlegi rendszerek gyakran nem nyújtanak megfelelő támogatást a receptek kezelésére, megosztására és értékelésére, különösen akkor, ha figyelembe vesszük az allergén információk fontosságát. Egy jól megtervezett alkalmazás nemcsak megkönnyítené a receptek megosztását és kezelését, hanem allergén információk megjelenítésével biztonságosabbá is tenné a főzést azok számára, akik ételallergiával élnek.

\section{A dolgozat célja}

A dolgozat célja egy olyan könnyen kezelhető webes alkalmazás, amely segítségével a felhasználók egyszerűen feltölthetik, megoszthatják és szerkeszthetik saját receptjeiket. Az alkalmazásnak támogatnia kell a receptek keresését és szűrését, különösen az allergén információk alapján, hogy az ételallergiával élő felhasználók is biztonságosan használhassák a platformot. Ezen szűrők úgy legyenek kialakítva, hogy azok kövessék a már bevett gyakarolatokat a vendéglátási iparban, ahol a ételallergének jelölése kiemelt fontosságú\cite{eu1169}.

Enek kívül egy olyan modern megoldásokat alkalmazó rendszer létrehozása, amely könnyen karbantartható és bővíthető, valamint megfelel a jelenlegi webes fejlesztési szabályoknak és bevett gyakorlatoknak. Emellett fontosnak tartom azt, hogy a rendszer könnyedén elindítható legyen helyi környezetben, és ami minimalizálja azon külső szoftvereket ami szükséges a rendszer futtatásához.
